\section{A Mercy Shown to the Saved}\label{sec:purgatory}

Between the particular judgment and the vision of God there stands, for many of the saved, a purification. The Church has taught it at two councils in nearly the same words. At Lyons in 1274 and again at Florence in 1439 she confessed that those who die truly penitent in charity, before they have made satisfaction by worthy fruits of penance for what they did and left undone, are purified after death by purgatorial pains; Florence's words are \latin{eorum animas poenis purgatoriis post mortem purgari}, where the Lyons profession had added after \latin{purgatoriis} the Greek-derived synonym \latin{seu catharteriis} and a clause naming the friar who had explained it.\endnote{Second Council of Lyons, profession of Michael Palaeologus, DS 856; Council of Florence, bull \emph{Laetentur caeli} (6 July 1439), DS 1304.} The Council of Trent, in its last session, recalled that the Church, ``instructed by the Holy Ghost, has, from the sacred writings, and the ancient tradition of the Fathers, taught, in sacred Councils, and very recently in this oecumenical Synod, that there is a Purgatory, and that the souls there detained are helped by the suffrages of the faithful, but principally by the acceptable sacrifice of the altar.''\endnote{Council of Trent, session~XXV (3--4 December 1563), \emph{Decree concerning Purgatory}, trans.\ J.~Waterworth, \emph{The Canons and Decrees of the Sacred and Oecumenical Council of Trent} (London, 1848; Burns and Oates reprint), pp.~232--233; DS 1820.}

Where purgatory stands matters more than anything else about it. In \latin{Benedictus Deus} it appears inside the definition of the blessed: the souls who had something to be purged are in heaven \latin{cum post mortem suam fuerint purgatae}, when they have been purged after their death. The Catechism puts it without qualification: ``All who die in God's grace and friendship, but still imperfectly purified, are indeed assured of their eternal salvation; but after death they undergo purification, so as to achieve the holiness necessary to enter the joy of heaven. The Church gives the name Purgatory to this final purification of the elect, which is entirely different from the punishment of the damned.''\endnote{\emph{Catechism of the Catholic Church} 1030--1031, official English.}

Purgatory, then, is not a second chance, a probation, or an appeal. No one enters it whose case is still open, and no one is ever lost from it. It lies wholly on the far side of the particular judgment, as the working out of a sentence already passed in the soul's favour. The tradition kept the mercy of purgatory by making it a mercy shown to the saved. A soul there has heard the word ``Come, ye blessed of my Father''; it is on its way.

\subsection{The cleansing fire}

The Church's tradition, the Catechism says, ``by reference to certain texts of Scripture, speaks of a cleansing fire.''\endnote{\emph{Catechism of the Catholic Church} 1031.} The first of those texts is Paul's word on the foundation and what is built on it:

\begin{quote}
Every man's work shall be manifest. For the day of the Lord shall declare it, because it shall be revealed in fire. And the fire shall try every man's work, of what sort it is. If any man's work abide, which he hath built thereupon, he shall receive a reward. If any man's work burn, he shall suffer loss: but he himself shall be saved, yet so as by fire. (1~Cor 3:13--15)
\end{quote}

Gregory the Great allows it to be read of the fire to come. ``But yet we must believe that before the day of judgment there is a Purgatory fire for certain small sins,'' he writes in the \emph{Dialogues}, and he reasons from Christ's saying that the sin against the Holy Ghost will not be forgiven ``neither in this world, nor in the world to come'': ``Out of which sentence we learn, that some sins are forgiven in this world, and some other may be pardoned in the next.''\endnote{Gregory the Great, \emph{Dialogues} IV.39 (trans.\ P.~W., ed.\ Gardner, 1911), p.~233.} He names the kind of faults he means: ``daily idle talk, immoderate laughter, negligence in the care of our family \ldots\ ignorant errors in matters of no great weight.'' And he adds a warning that turns the doctrine back upon the living: ``none can be there purged, no, not for the least sins that be, unless in his lifetime he deserved by virtuous works to find such favour in that place.''\endnote{Gregory, \emph{Dialogues} IV.39 (1911), pp.~233--234.} The Catechism takes its own words on the purifying fire from this chapter.\endnote{\emph{Catechism of the Catholic Church} 1031, citing Gregory, \emph{Dialogues} 4, 39.}

Augustine had taught the same truth plainly before him. Against those who would make every punishment after death purgatorial, he writes in \emph{The City of God}: ``temporary punishments are suffered by some in this life only, by others after death, by others both now and then; but all of them before that last and strictest judgment.''\endnote{Augustine, \emph{De civitate Dei} XXI.13 (Dods).} His reserve concerned only the reading of Paul's verse. He referred the fire of 1~Corinthians~3 chiefly to the tribulations of this life, which try the loves a Christian has built on Christ and burn away whatever was earthly in them; and when he came to the opinion that the spirits of the departed pass through such a fire in the interval between death and the resurrection, he answered with a reverent reserve that the Latin text preserves exactly: \latin{non redarguo, quia forsitan uerum est}---``this I do not contradict, because possibly it is true.''\endnote{Augustine, \emph{De civitate Dei} XXI.26; Latin from the edition of Emanuel Hoffmann, \emph{Corpus Scriptorum Ecclesiasticorum Latinorum} 40.2 (Vienna, 1900); English from Dods (1871).} The Church has taught in council what Augustine already held: there is a purification after death, and the prayers of the living help those who undergo it.

The Church leaves open what this fire is, and her teachers have spoken of it in more than one way. Gregory speaks of it simply as fire, and the Latin schools after him held it to be a real fire: the appendix to the Supplement of the \emph{Summa} names the pain of sense in purgatory ``punishment by corporeal fire.''\endnote{\emph{Summa theologiae}, Supplement, Appendix~I, q.~2, a.~1, corpus (English Dominican translation, 1922).} The Council of Florence, defining the doctrine for Latins and Greeks together, speaks only of purgatorial pains and does not name the fire. Benedict~XVI records the view of ``some recent theologians'' that ``the fire which both burns and saves is Christ himself, the Judge and Saviour,'' and he develops it with great beauty: ``Before his gaze all falsehood melts away. This encounter with him, as it burns us, transforms and frees us, allowing us to become truly ourselves. \ldots\ But it is a blessed pain, in which the holy power of his love sears through us like a flame, enabling us to become totally ourselves and thus totally of God.''\endnote{\emph{Spe salvi} 47, official English.}

\subsection{Joy in the fire}

No saint has spoken of purgatory from within more movingly than Catherine of Genoa. Her \emph{Treatise on Purgatory} rests on what she says God made known to her ``by special grace''; it is the private revelation of a canonized saint, which the Church honours but has not made her own doctrine. It describes the souls there as wholly given over to the will of God, unable any longer to sin or to wish anything but what God wills. Her teaching holds together two things that seem incompatible. The souls suffer, and they rejoice:

\begin{quote}
I do not believe it would be possible to find any joy comparable to that of a soul in purgatory, except the joy of the blessed in paradise,---a joy which goes on increasing day by day, as God more and more flows in upon the soul, which He does abundantly in proportion as every hindrance to His entrance is consumed away. The hindrance is the rust of sin; the fire consumes the rust, and thus the soul goes on laying itself open to the Divine inflowing.\endnote{Catherine of Genoa, \emph{The Treatise on Purgatory}, chapter~2, translated from the Italian, with a preface by H.~E. Manning (London: Burns and Lambert, 1858), p.~3.}
\end{quote}

And yet, she continues, ``they suffer a torment so extreme, that no tongue could describe it, no intellect could form the least idea of it, if God had not made it known by special grace.''\endnote{Catherine of Genoa, \emph{Treatise on Purgatory}, chapter~2 (1858), p.~4.} The pain is the pain of love delayed: the soul sees God as its whole good and is held back for a time from him by what remains in it of sin. The joy is the joy of love that knows it is going home. Both grow from the same root, and neither cancels the other.

Here is the deepest reason why purgatory is a mercy. No one would wish to enter the presence of God carrying what he would be ashamed to bring there. The \emph{Imitation of Christ} draws the practical conclusion for the living: ``Better is it now to purify the soul from sin, than to cling to sins from which we must be purged hereafter.''\endnote{Thomas \`a Kempis, \emph{The Imitation of Christ}, book~I, chapter~24, 2, trans.\ William Benham (Project Gutenberg eBook 1653).} The same purification can be begun here, by penance patiently borne, by injuries forgiven, by daily faults confessed and amended. Every suffering accepted in union with Christ is purgatory begun on this side of death.
